\documentclass[a4paper, 12pt, UTF8]{ctexart}
\usepackage{mathtools}
\usepackage{amssymb}
\usepackage{amsthm}
\usepackage[margin = 2.5cm]{geometry}
\usepackage{tikz}
\usepackage{graphicx}
\usepackage{xcolor}

\newtheoremstyle{theoremstyle}
  {}{}{\itshape\color{red}}{}{\bfseries}{.}{\newline}{}
\theoremstyle{theoremstyle}
\newtheorem{theorem}{theorem}

\newtheoremstyle{propertystyle}
  {}{}{\itshape\color{magenta}}{}{\bfseries}{.}{\newline}{}
\theoremstyle{propertystyle}
\newtheorem{property}[theorem]{property}

\newtheoremstyle{definitionstyle}
  {}{}{\normalfont\color{teal}}{}{\bfseries}{.}{\newline}{}
\theoremstyle{definitionstyle}
\newtheorem{definition}[theorem]{definition}

\newtheoremstyle{algorithmstyle}
  {}{}{\ttfamily\color{blue}}{}{\bfseries}{.}{\newline}{}
\theoremstyle{algorithmstyle}
\newtheorem{algorithm}[theorem]{algorithm}



\title{transformer}
\author{Yifei He}
\date{\today}

\begin{document}
\maketitle
首先提到翻译问题 要求神经网络的输入维度和输出维度是任意的。
首先让一个模型可以处理任意维度的输入，token化。
不同的信息之间肯定需要信息传递，但是单纯通过残差连接值矩阵 模型就太简单了
我们需要引入权重，权重应该和信息传递的输出和接受方有关 同时应该是一个实数。于是我们想到kv矩阵 点积构造权重
为了让单个token内部传递信息 只需要mlp
不同token之间传递信息需要位置无关。

该架构应该无法使用fnn表示

于是我们得到了一个输入任意维度 输出为对应维度的nn
我们事先指定小于输入维度个数的输出为我们需要的 就得到了一个输入为任意维度 输出为指定维度的nn
很容易的 该模型非常适合用于预测，
通过自回归，滚动预测任意长度的值。
如果我们让模型自己决定何时停止输出，这就是一个无输入，输出模型自己决定的维度的nn
自然语言输出也类似于预测 所以可以用这样的模型表示。

回到最初的问题，我们需要的模型是输入任意维度 输出任意维度, 所以我们还需要在上述模型的基础上添加一个输入。
这就是cross attention部分。我们可以将翻译任务分为两部分，
首先是理解原文，转化为一个表示其含义的向量，
然后是根据目标语言的语法知识来输出译文。
我们之前提到的decoder only transformer 就可以用于第二步 且输出可以是任意维度。
所以我们只需要类似self attention 但将信息的发射端改为完全不同的token空间 就可以实现任意维度信息的嵌入

\bibliographystyle{plain}
\bibliography{refs}
\end{document}